\section{Conclusion and Future Work}

\subsection{What was built}

This project set out to answer a practical question: can a constrained
combinatorial encoding of graph pruning, combined with trust-aware learning
on the surviving topology, improve task allocation in a dynamic rover
network? To answer it honestly we built the complete pipeline rather than a
toy demonstration. The system consists of a constrained QUBO formulation of
edge selection with degree, budget, flow connectivity, and flow capacity
constraints; a feasibility-projected production annealer that moves only in
the space of valid edge sets; a trust-aware graph neural network that scores
rovers on the pruned topology; a capacity-aware greedy assignment stage; and
an online trust model that folds each mission's outcomes back into the next
cycle's features. Around this core we built the apparatus needed to audit
it: a causal simulation harness in which mission success genuinely depends
on the kept topology, an exact MILP baseline for ground truth, a hostile
test suite of 491 checks organised into fourteen sections, a physics
integration based on Gazebo Harmonic and ROS 2 Jazzy, and an interactive
browser console driven by a custom three.js simulation engine that runs the
same allocation logic against a live visualisation.

\subsection{What was proven}

The QUBO encoding itself is verified three independent ways: an independent
re-derivation of the energy function agrees with the constructed $Q$ matrix
term by term; brute-force enumeration of all edge subsets on small instances
reproduces the encoding's optima; and an exact linearized solve of the
matrix with CBC confirms the same energies. The production annealer is 100\%
exact on small instances against enumeration, 95\% exact against MILP on
twenty larger instances with a mean true optimality gap of $7\times10^{-5}$,
and 87\% exact across all sensitivity perturbations of temperatures,
cooling, restarts, penalty scales, and move probabilities. At the system
level, trust-aware assignment beats random assignment by $+13.8$ percentage
points of mission success and a trust-only heuristic by $+8.13$ points, with
a paired bootstrap 95\% confidence interval of $[+4.2, +11.7]$ and a
permutation $p$-value of $0.016$. Trust tracking itself is worth $+13.75$
points of staleness margin: the benefit persists when topology decisions are
amortised over several cycles, which matters when the solve must be
offloaded to a cloud tier.

\subsection{What was disproven or humbled}

An honest audit must report what did not work, and this project treats those
results as first-class findings. First, classical MILP solved by CBC is
exact and roughly one hundred times faster than the annealer on the same
instances. The QUBO formulation is therefore a hardware-target encoding, not
the recommended classical solver, and the paper says so without
qualification. Second, four serious annealing samplers, a bit-level
simulated annealer, a production simulated annealing library, a tabu search,
and simulated quantum annealing, reached the optimum on the penalty-encoded
landscape only 48\% to 64\% of the time, and random bit-string restarts
returned infeasible states in 19 of 25 instances. Third, at the ten-rover
scale a simple threshold pruner beats the QUBO pruner on raw mission
success, 38.2\% to 33.8\% with $p = 0.028$, because the threshold keeps
redundant links that rescue delivery through bottleneck topologies. Fourth,
message passing adds only $+1.7$ insignificant points over a no-graph
readout, so the graph mechanism's value is not yet established at this
scale. Fifth, a cloud-trained readout is 3.9 points worse than the fixed
hand-designed weights ($p = 0.49$), a clean null result against the
trainability hypothesis.

\begin{figure}[htbp]
\centering
\begin{tikzpicture}[
  font=\small,
  col/.style={draw, rounded corners=2pt, inner sep=8pt, align=left,
              text width=6.4cm},
]
\node[col] (left) {%
  \textbf{What we set out to prove}\\[2pt]
  $\bullet$ Constrained QUBO pruning is a viable edge-selection encoding\\
  $\bullet$ Annealing samplers solve the encoded problem reliably\\
  $\bullet$ Graph message passing beats a flat readout\\
  $\bullet$ A trained readout beats fixed weights};
\node[col, right=1.2cm of left] (right) {%
  \textbf{What we actually proved}\\[2pt]
  $\bullet$ Encoding verified three ways; projected annealer 95 to 100\%
    exact\\
  $\bullet$ Generic samplers reach the optimum only 48 to 64\% of the time\\
  $\bullet$ Message passing adds $+1.7$ points, not significant\\
  $\bullet$ Trained readout is 3.9 points worse ($p = 0.49$)};
\draw[-{Stealth[length=3mm]}, thick] (left.east) to (right.west)
  node[midway, above, font=\footnotesize\itshape] {audit};
\end{tikzpicture}
\caption{The distance between the proposal and the audited result. Every
item on the left survived contact with measurement in a modified form; the
right column is what the evidence supports.}
\label{fig:conclusion-audit}
\end{figure}

\subsection{Lessons}

Four lessons generalise beyond this project. First, verify encodings
against independent ground truth before trusting them: our encoding survived
because it was cross-checked against re-derivation, enumeration, and an
exact solver, and that triple verification is what lets the downstream
claims stand. Second, penalty-encoded constrained QUBOs are a hard landscape
for generic samplers; the penalty barriers that enforce feasibility are
exactly what trap samplers that move in raw bit space. Third,
feasibility-aware move sets matter more than solver sophistication: the
projected annealer's advantage came from proposing only valid moves and
repairing candidates, not from a cleverer temperature schedule. Fourth,
negative results deserve the same prominence as positive ones, because the
negative results here, the MILP dominance, the threshold pruner's win, and
the readout-training null, are the findings most likely to save a future
team from repeating the same detours.

\subsection{Future work}

Several directions follow directly from what the audit exposed. The Gazebo
integration should publish velocity commands so that physics drives the
planning inputs rather than replaying scripted poses, closing the loop
between the simulator and the allocator. The browser console needs WebSocket
streaming in place of request-response stepping, plus session persistence so
long missions survive reconnects. The pruning formulation should be extended
to hierarchical or clustered variants, since connectivity is currently
enforced only up to twenty rovers due to the dense memory footprint of the
matrix. The $Q$ matrix should be run on real annealing hardware from D-Wave
to test whether the hardware changes the sampler-reliability picture that
the classical samplers painted. The readout-training null result deserves a
revisit with larger and more diverse training data, since the current
failure may be a data problem rather than a model problem. Finally, the
deterministic link-quality model should be replaced with richer radio models
that include fading and interference, which would stress the trust model in
ways the current environment cannot.

\subsection{Closing remarks}

Building this system taught the team that the distance between a plausible
pipeline and a trustworthy one is measured in verification infrastructure,
not in model complexity. The parts of the project we are most confident in
are not the cleverest components; they are the components pinned down by
independent ground truth, and the parts we trusted too early were the ones
that measurement later humbled. We set out to show that a quantum-inspired
encoding could earn its place in a rover network, and we ended up with
something more useful: a complete, audited picture of where the encoding
helps, where it does not, and exactly how much each design decision is
worth. That picture, including its negative regions, is the contribution.

\subsection{Pending items in this report}

The following placeholders remain in this report and must be completed before submission:

\begin{itemize}
  \item title page author names and date;
  \item the abstract text;
  \item Gazebo world and status-message screenshots;
  \item fleet console screenshots (orbit view, mission strip), performance capture, and demo video link;
  \item the reference list for the final section (to be supplied separately).
\end{itemize}
